\section*{\S~9 Kampfrollen}
\label{sec:kampfrollen}
\addcontentsline{toc}{section}{\S~9 Kampfrollen}

\begin{enumerate}[label=(\arabic*)]
    \item Kampfrollen stellen die Ausbildungen für verschiedene Funktionen, die in einer Mission ausgeübt werden können, dar. Welche dies sind, wird durch die Leitung festgelegt. Die Kampfrollen Schütze (Grundausbildung), Gruppenführer und Missionsbauer/Zeus sind gesetzt.
    \begin{enumerate}[label=(\Alph*)]
        \item Sonder-Kampfrollen sind Kampfrollen, für die neben der Schützen-Kampfrolle noch weitere Ausbildungen benötigt werden. Diese Voraussetzungen legt die Leitung zusammen mit den Ausbildern für diese Sonder-Kampfrollen fest. Die Voraussetzungen müssen klar und in geeigneter Weise mit allen Mitgliedern kommuniziert werden.
    \end{enumerate}
    \item Alle Kampfrollen benötigen einen Ausbildungsleitfaden, welcher als stichpunktartige Checkliste der Ausbildungsinhalte zu verstehen ist. Neben den Bedingungen in Arma 3 müssen die Aufgaben, die allgemeinen Verhaltensweisen, das Verhalten in der Gruppe und das Verhalten im Feuerkampf Teil dieses Ausbildungsleitfadens sein. Besonders komplizierte Kampfrollen können um ein Handbuch für Ausbilder und Inhaber ergänzt werden. Ob dies nötig ist, legen die Leitung oder die bestehenden Ausbilder fest.
    \begin{enumerate}[label=(\Alph*)]
        \item Sonder-Kampfrollen benötigen ausschließlich ein Handbuch.
        \item Bei großen Änderungen an einem Ausbildungsleitfaden/Handbuch können die Ausbilder die Eigentümer einer Kampfrolle zu einer Nachschulung verpflichten. Nimmt ein Mitglied nicht daran teil, kann der Stab diesem die Kampfrolle aberkennen.
        \item Mitglieder können selbständig mit Vorschlägen für Änderungen, oder Neueinführungen für Kampfrollen an die Leitung herantreten.
        \item Die Leitung hat das Recht, die Erstellung oder Überarbeitung eines Ausbildungsleitfaden/Handbuchs an die Ausbilder oder Interessierte zu delegieren.
        \item Ausbildungsleitfäden/Handbücher müssen durch die Leitung, oder ein anderes, von der Leitung damit beauftragtes Gremium geprüft und beschlossen werden.
        \item Die Leitung zeichnet sich für die Ausbildungsleitfäden/Handbücher verantwortlich.
    \end{enumerate}
    \item Jedes Mitglied, dass die Voraussetzungen für den Dienstgrad Gefreiter erfüllt hat, auch wenn die Beförderung noch nicht stattgefunden hat, hat ein Recht auf Ausbildung, sofern es die Voraussetzungen für diese Kampfrolle erfüllt.
    \item In wie weit eine Kampfrolle durch eine Überprüfung, egal welcher Art, bestätigt werden muss, legt die Leitung fest. Sonder-Kampfrollen benötigen einer Überprüfung.
    \item Jedes Mitglied, welches eine Kampfrolle erworben hat, kann gegen ein anderes Mitglied, welches die gleiche Kampfrolle erworben hat, eine Kampfrollen bezogene Beschwerde, wegen Fehlerverhalten in dieser Kampfrolle, einreichen. In einem solchen Fall tritt der Stab gemeinsam mit den Ausbildern für diese Kampfrolle zusammen und beurteilt diesen Fall. Wird ein Fehlverhalten festgestellt, kann je nach Schwere des Fehlverhaltens eine Nachschulung für diese Kampfrolle, die Sperrung eines Mitglieds für diese Kampfrolle, oder die Aberkennung dieser Kampfrolle erfolgen. Der Stab trifft hierbei die Entscheidung. Die Ausbilder der Kampfrolle haben lediglich eine beratende Funktion.
    \begin{enumerate}[label=(\Alph*)]
        \item Wenn das Fehlerverhalten einen schweren Regelverstoß nach \S~8 Absatz \ref{itm:schwerer_verstoß} darstellt, kann zusätzlich eine Bestrafung nach \hyperref[sec:umgang_mit_regelverstoeßen]{\S~8} erfolgen.
        \item Richtet sich die Beschwerde gegen einen Ausbilder dieser Kampfrolle, so kann dieser auch seinen Posten enthoben werden.
    \end{enumerate}
\end{enumerate}